\documentclass{beamer}
\usepackage[utf8]{inputenc}

\usetheme{Madrid}
\usecolortheme{default}
\usepackage{amsmath,amssymb,amsfonts,amsthm}
\usepackage{txfonts}
\usepackage{tkz-euclide}
\usepackage{listings}
\usepackage{adjustbox}
\usepackage{array}
\usepackage{tabularx}
\usepackage{gvv}
\usepackage{lmodern}
\usepackage{circuitikz}
\usepackage{tikz}
\usepackage{graphicx}
\usepackage{multicol}

\setbeamertemplate{page number in head/foot}[totalframenumber]

\usepackage{tcolorbox}
\tcbuselibrary{minted,breakable,xparse,skins}



\definecolor{bg}{gray}{0.95}
\DeclareTCBListing{mintedbox}{O{}m!O{}}{%
  breakable=true,
  listing engine=minted,
  listing only,
  minted language=#2,
  minted style=default,
  minted options={%
    linenos,
    gobble=0,
    breaklines=true,
    breakafter=,,
    fontsize=\small,
    numbersep=8pt,
    #1},
  boxsep=0pt,
  left skip=0pt,
  right skip=0pt,
  left=25pt,
  right=0pt,
  top=3pt,
  bottom=3pt,
  arc=5pt,
  leftrule=0pt,
  rightrule=0pt,
  bottomrule=2pt,
  toprule=2pt,
  colback=bg,
  colframe=orange!70,
  enhanced,
  overlay={%
    \begin{tcbclipinterior}
    \fill[orange!20!white] (frame.south west) rectangle ([xshift=20pt]frame.north west);
    \end{tcbclipinterior}},
  #3,
}
\lstset{
    language=C,
    basicstyle=\ttfamily\small,
    keywordstyle=\color{blue},
    stringstyle=\color{orange},
    commentstyle=\color{green!60!black},
    numbers=left,
    numberstyle=\tiny\color{gray},
    breaklines=true,
    showstringspaces=false,
}
%------------------------------------------------------------
%This block of code defines the information to appear in the
%Title page
\title %optional
{7.4.36}
\date{October 12,2025}


\author 
{Jnanesh Sathisha karmar - EE25BTECH11029}



\begin{document}



\frame{\titlepage}
\begin{frame}{Question:}


\noindent The abscissa of two points \textbf{A} and \textbf{B} are roots of the equation $x^2 + 2ax - b^2 = 0$ and their ordinates are roots of the equation $x^2 + 2px - q^2 = 0$. Find the equation and the radius of the circle with $AB$ as diameter.
\end{frame}

\begin{frame}{Theoretical Solution}

 Define the Problem in Vector Form

The vector equation for a circle is given by:
\begin{align*}
    \norm{\vec{x}}^2 + 2\vec{u}^T \vec{x} + f = 0
\end{align*}
where $\vec{x} = \myvec{x \\ y}$, the center is $\vec{c} = -\vec{u}$, and $f = \Vert \vec{u} \Vert^2 - r^2$.
Let the points be $\vec{x}_A = \myvec{x_1 \\ y_1}$ and $\vec{x}_B = \myvec{x_2 \\ y_2}$.

 Determine the Vector $\vec{u}$

Since $AB$ is the diameter, the center $\vec{c}$ is the midpoint of $\vec{x}_A$ and $\vec{x}_B$.
\begin{align*}
    \vec{c} = \frac{\vec{x}_A + \vec{x}_B}{2} = \frac{1}{2} \myvec{x_1 + x_2 \\ y_1 + y_2}
\end{align*}
\end{frame}
\begin{frame}{Theoretical Solution}
Using Vieta's formulas for the given quadratic equations:
\begin{itemize}
    \item For $x^2 + 2ax - b^2 = 0$:
        \begin{align*}
            x_1 + x_2 &= -2a \\
            x_1 x_2 &= -b^2
        \end{align*}
    \item For $y^2 + 2py - q^2 = 0$:
        \begin{align*}
            y_1 + y_2 &= -2p \\
            y_1 y_2 &= -q^2
        \end{align*}
\end{itemize}
Substituting the sum of roots into the center vector $\vec{c}$:
\begin{align*}
    \vec{c} = \frac{1}{2} \myvec{-2a \\ -2p} = \myvec{-a \\ -p}
\end{align*}
Therefore, the vector $\vec{u}$ is:
\begin{align*}
    \vec{u} = -\vec{c} = - \myvec{-a \\ -p} = \myvec{a \\ p}
\end{align*}

\end{frame}
\begin{frame}{Theoretical Solution}
 Determine the Radius $r$ and Scalar $f$

The squared diameter $d^2$ is $\norm{\vec{x}_A - \vec{x}_B}^2$.
\begin{align*}
    d^2 = (x_1 - x_2)^2 + (y_1 - y_2)^2
\end{align*}
We use the identity $(m-n)^2 = (m+n)^2 - 4mn$.
\begin{align*}
    (x_1 - x_2)^2 &= (x_1 + x_2)^2 - 4(x_1 x_2) \\
    &= (-2a)^2 - 4(-b^2) \\
    &= 4a^2 + 4b^2 = 4\brak{a^2 + b^2}
\end{align*}
\begin{align*}
    (y_1 - y_2)^2 &= (y_1 + y_2)^2 - 4(y_1 y_2) \\
    &= (-2p)^2 - 4(-q^2) \\
    &= 4p^2 + 4q^2 = 4\brak{p^2 + q^2}
\end{align*}
The squared diameter is:
\begin{align*}
    d^2 = 4\brak{a^2 + b^2} + 4\brak{p^2 + q^2} = 4\brak{a^2 + b^2 + p^2 + q^2}
\end{align*}

\end{frame}
\begin{frame}{Theoretical Solution}
The squared radius $r^2$ is $d^2/4$:
\begin{align*}
    r^2 = a^2 + b^2 + p^2 + q^2
\end{align*}
Now, we find the scalar $f = \norm{\vec{u}} - r^2$.
\begin{align*}
    \norm{\vec{u}} = \norm{ \myvec{a \\ p}}^2 = a^2 + p^2
\end{align*}
\begin{align*}
    f &= \brak{a^2 + p^2} - \brak{a^2 + b^2 + p^2 + q^2} \\
    &= a^2 + p^2 - a^2 - b^2 - p^2 - q^2 \\
    &= -(b^2 + q^2)
\end{align*}


\end{frame}
\begin{frame}{Theoretical Solution}
Final Equation and Radius
Substitute $\vec{u}$ and $f$ into the circle's vector equation:
\begin{align*}
    \norm{\vec{x}}^2 + 2\vec{u}^T \vec{x} + f = 0 \\
    \norm{\myvec{x \\ y}}^2 + 2 \myvec{ a & p } \myvec{x \\ y} - \brak{b^2 + q^2} = 0 \\
    \brak{x^2 + y^2} + 2\brak{ax + py} - \brak{b^2 + q^2} = 0 \\
    x^2 + y^2 + 2ax + 2py - b^2 - q^2 = 0
\end{align*}
\textbf{Equation of the circle:}
\begin{align*}
    x^2 + y^2 + 2ax + 2py - b^2 - q^2 = 0
\end{align*}
\textbf{Radius of the circle:}
\begin{align*}
    r = \sqrt{a^2 + b^2 + p^2 + q^2}
\end{align*}
    
    
\end{frame}

\begin{frame}[fragile]
    \frametitle{python Code }

    \begin{lstlisting}
import sympy

def solve_circle_problem():
    x, y, a, b, p, q = sympy.symbols('x y a b p q')
    x1, x2, y1, y2 = sympy.symbols('x1 x2 y1 y2')

    print("--- Problem Setup ---")
    print("Equation for x-coordinates: x**2 + 2*a*x - b**2 = 0")
    print("Equation for y-coordinates: y**2 + 2*p*y - q**2 = 0\n")

    sum_x = -2 * a
    prod_x = -b**2

    sum_y = -2 * p
    prod_y = -q**2

    print("--- Applying Vieta's Formulas ---")
    print(f"Sum of x-roots (x1 + x2): {sum_x}")
    print(f"Sum of y-roots (y1 + y2): {sum_y}\n")

   
   


    \end{lstlisting}
\end{frame}

\begin{frame}[fragile]
    \frametitle{python Code }
    \begin{lstlisting}
     center_h = sympy.simplify(sum_x / 2)
    center_k = sympy.simplify(sum_y / 2)
    print(f"--- Circle Center (h, k) ---")
    print(f"h = (x1 + x2) / 2 = {center_h}")
    print(f"k = (y1 + y2) / 2 = {center_k}\n")
    
    x_diff_sq = sum_x**2 - 4 * prod_x
    y_diff_sq = sum_y**2 - 4 * prod_y
    
    diameter_sq = x_diff_sq + y_diff_sq
    radius_sq = sympy.simplify(diameter_sq / 4)
    radius = sympy.sqrt(radius_sq)

    print("--- Circle Radius (r) ---")
    print(f"Squared diameter d^2 = (x1-x2)^2 + (y1-y2)^2 = {diameter_sq}")
    print(f"Squared radius r^2 = d^2 / 4 = {radius_sq}")
    print(f"Radius r = {radius}\n")

\end{lstlisting}
\end{frame}
\begin{frame}[fragile]
    \frametitle{python Code }
    \begin{lstlisting}
   circle_eq_standard = (x - center_h)**2 + (y - center_k)**2 - radius_sq
    
    circle_eq_general = sympy.expand(circle_eq_standard)

    print("--- Final Equation of the Circle ---")
    print(f"General Form: {circle_eq_general} = 0")


if __name__ == "__main__":
    solve_circle_problem()

\end{lstlisting}
\end{frame}


\end{document}
